\chapter{Einleitung}

\thispagestyle{standard}
\pagestyle{standard}
\section{Motivation}

Die zunehmende Vernetzung und die immer stärkere Integration von Systemen macht
auch vor dem Bereich der Automatisierungstechnik nicht Halt. So ist hier in den letzten
Jahren ein steigender Bedarf an Interkonnektivität der einzelnen, vormals autonomen
Systeme zu verzeichnen. Vorschub leistet dieser Bewegung, dass durch die gestiegene
Verfügbarkeit von breitbandigen Verbindungen auch der Wunsch nach höherer Integrationsdichte
gewachsen ist. Immobilien, Ladengeschäfte, Produktionsstätten oder ganze
Firmenstandorte werden zunehmend als Verwaltungseinheiten gesehen, für die eine Administration
und Überwachung aus der Ferne oder die Einbindung in Firmennetze erfolgen
soll.
Auch der immer häufigereWunsch, Geschäftsprozesse automatisiert in Software abzubilden,
lässt es sinnvoll erscheinen, auf verschiedene Systeme eines Standorts aus der Ferne
als Entität zugreifen zu können. Damit verändern sich auch die Anforderungen an die
Schnittstellen eines Systems. Waren früher systemindividuelle Bediener vor Ort für die
Betreuung ausreichend, so soll dies heute immer öfter auch aus der Ferne zentral über
eine standardisierte Schnittstelle möglich sein. Somit hält die Problematik der verteilten
Anwendungen auch in vormals oft vollständig unabhängigen Systemen Einzug.
Gerade in der Automatisierungstechnik macht sich bemerkbar, dass der scheinbare Wildwuchs
an konkurrierenden Bus Systemen der letzten Jahre zu einer extremen Inkompatibilität
der Systeme geführt hat. Systeme die nun wieder eine Einheit im Sinne der remoten
Verwaltung bilden sollen. An Hand der Energieverwaltung eines Gebäudes oder
einer Produktionsstätte lässt sich dies veranschaulichen. Für eine zentrale Überwachung



\section{Zielsetzung}

\section{Anmerkung}

\chapter{Grundlagen}

\section{OpenNES}

\subsection{Architektur}

\subsection{}

\section{ModBus Protokoll}

\section{Abbildungen}

In Moby-Dick geht es in erster Linie um die Jagd auf einen weißen Pottwal \footnote{nach: https://de.wikipedia.org/wiki/Moby-Dick}. Kurze Zitate (unter drei Zeilen) müssen mit Anführungszeichen gekennzeichnet werden. Außerdem müssen die Quelle sowie die Seite angegeben werden. Ein Beispiel für ein kurzes Zitat: \glqq Komisch. Manch einer von uns wünschte sich, er lebe auf einer Südseeinsel.\grqq \cite{MELVILLE:MOBYDICK1997} (Seite 100).
\section{Leck mich}
  \begin{quote}
"Das Buch hier Lieblingsbuch. Viele Blätter viele, schöne Bilder. Du kennen Worte?"\newline
"Ja"\newline
"Ich kennen Bilder. Das ein Wal. Du lesen Worte!"\newline
"Durch das Herz des Wals strömt mehr Flüssigkeit als durch das große Wasserleitungsrohr unter der London Bridge, jedoch strömt das Wasser nicht so stark, wie das Blut, das vom Herz des Wals pocht."\newline
"Du gut, ich danken dir."  \upshape \cite{MELVILLE:MOBYDICK1997} (Seite 500)
  \end{quote}

Ein sehr praktisches Package ist cleveref. Es automatisiert und erleichtert das setzen von Referenzen ungemein. Als Beispiel wird eine Referenz auf das FHS-Logo gesetzt siehe \cref{FIG_LOGO}.

Für das Verfassen von wissenschaftlichen Arbeiten können eine Vielzahl an Quellentypen herangezogen werden. Beispiele hierfür sind Leitfäden (Manuals) \citep{RFC2828} \citep{80211i} \citep{80211} \cite{X800} \citep{TR102377} \citep{EN301893} \citep{PUB197} \citep{PUB74}, Bücher \citep{Fis04a} \citep{Rei05a} \citep{Tan00a} \citep{Ste04a} \citep{GMS00a} \citep{HL98a}, Sammelbände \citep{EHL00a} \citep{Sch94a}, Journal-Artikel \citep{TM03a} \citep{CP03a}, Konferenz-Proceedings \citep{HCB00a} \citep{KBW04a} \citep{KSW04a} \citep{HK05a}, Internetmagazine \citep{Eke05a}, Webquellen \cite{nist} \cite{php} \cite{BDKMT93a} \cite{IDSSM} sowie Diplomarbeiten und Dissertationen \citep{Sch98} \cite{Hae94a}.

Als Beispiel für eine Abkürzung wird hier \ac{ADF} angeführt. Das Package schreibt automatisch das erste Vorkommen der Abkürzung aus. Die zweite Verwendung von \ac{ADF} wird also abgekürzt. Ist ein Ausschreiben einer Abkürzung gewünscht wird der acl-Befehl verwendet. Dies führt zu \acl{ADF}. Abkürzungen müssen in der Datei \glqq 05Abkuerzungsverzeichnis\grqq angegeben werden.

\section{Quelltext}

\texttt{printf("Hallo Welt")} für Ausschnitte von Sourcecode innerhalb von Text

\lstset{escapeinside={\%*}{*)},numbers=none}%oder numbers=left
\begin{lstlisting}[language=C,
caption=Beispiel-Listing,
label=LST_SAMPLE]
serverTCP = new TcpListener(IPAddress.Parse(serverIP), serverPort);
\end{lstlisting}

\lstset{escapeinside={\%*}{*)},numbers=left}
\lstinputlisting[language=Matlab, caption=Einfaches Matlabprogramm in einer Datei, label=list:hello.m]{Listings/hello.m}
\section{Text}

\section{Refik}

\section{Bilder}

\begin{figure}[H]
\begin{center}
	\includegraphics[scale=0.4]{BilderAllgemein/Logo.jpg}
\end{center}
	%\includegraphics[width=\textwidth]
	%\end{center}
	% Title
	\caption{Das FHS-Logo}
	% Unique name: identifier for referencing
	\label{FIG_LOGO}
\end{figure}

\section{Formeln}

Formeln sind für jeden Abschnitt rechtsbündig von dieser zu nummerieren, um einen späteren Bezug in der Arbeit zu gewährleisten. Formeln werden üblicherweise in "`Computer Modern Roman"' (\LaTeX{}-Standard) gesetzt. In diesem Template wird die Formel-Schrift bzw. das Package \texttt{eulervm} verwendet. Abgesetzte Formeln werden in \LaTeX{} durch die 
\emph{equation} Umgebung definiert. Formelausdrücke innerhalb von Textabschnitten erhält man durch \$Formel\$.

\subsection*{Beispiel}
%
Der \emph{Sinus cardinalis} oder sinc-Funktion ist eine mathematische Funktion $f$, welche in nicht-normierter Version als

\begin{equation}
	f(x) := \frac{\sin(x)}{x}
	\label{eq:bsp}
\end{equation}

definiert wird. In der digitalen Signalverarbeitung findet meistens nachfolgende normierte Version $\mathrm{si}(x)$ oder $\mathrm{sinc}(x)$ Anwendung \cite{x1}, \cite{x2}. Für eine Visualisierung dieser Funktionen siehe Abb.~\ref{FIG_LOGO}.

\begin{equation}
	f(x) := \frac{\sin(\pi x)}{\pi x}
	\label{EQ_SAMPLE}
\end{equation}

\section{Beispiel für Tabellen}
%
Es empfiehlt sich, für Tabellen die Standard-\LaTeX{}-Umgebung \emph{tabular} zu verwenden. Bei Bedarf können natürlich auch Erweiterungen (z.B.~\emph{tabularx} oder \emph{array}) zur Anwendung kommen. Eine mögliche Darstellung zeigt Tabelle \ref{Table_Sinc}.

\begin{table}[h!]%
	\begin{center}
	
		\begin{tabular}{|r|r|r|}
			\firsthline
			$x$&$\mathrm{sinc}(x)$&$\mathrm{sin}(x)$\\\hline\hline
			$-0.5$&0.6366&-0.4794\\\hline
			$0$&1.0000&0\\\hline
			$0.5$&0.6366&0.4794\\\hline
		\end{tabular}
		\caption{Zwei Werte der Sinc-Funktion}
		\label{Table_Sinc}
	\end{center}
\end{table}








